\documentclass[11pt]{article}
\usepackage[T1]{fontenc}
\usepackage[utf8]{inputenc}
\usepackage{lmodern}
\usepackage[margin=1in]{geometry}
\usepackage{amsmath,amssymb,amsthm}
\usepackage{booktabs,longtable,array}
\usepackage[hidelinks]{hyperref}
\usepackage{microtype}
\newtheorem{theorem}{Theorem}[section]
\newtheorem{lemma}[theorem]{Lemma}
\newtheorem{proposition}[theorem]{Proposition}
\theoremstyle{remark}
\newtheorem{remark}[theorem]{Remark}
\newcommand{\E}{\mathbb{E}}
\newcommand{\Prob}{\mathbb{P}}
\newcommand{\ind}{\mathbf{1}}
\newcommand{\Pois}{\operatorname{Poisson}}
\newcommand{\Bin}{\operatorname{Binomial}}
\newcommand{\Var}{\operatorname{Var}}
\newcommand{\en}{\mathcal{E}}
\newcommand{\code}[1]{\texttt{\small #1}}
\title{A Lean-verified upper bound of 4.268 for random 3-SAT}
\author{Fedor Vorobyev\\\small Independent researcher\\
  \small\href{mailto:fedorvorobev@gmail.com}{fedorvorobev@gmail.com}}
\date{October 2026}
\hypersetup{pdftitle={A Lean-verified upper bound of 4.268 for random 3-SAT},
  pdfauthor={Fedor Vorobyev}}
\begin{document}
\maketitle
\begin{abstract}
We prove that a uniformly random signed 3-CNF formula with
$\lfloor 4.268n\rfloor$ independent clauses on $n$ variables is
unsatisfiable with probability tending to one. Each clause contains three
distinct variables; repeated clauses are allowed. The proof specializes
the clause/site interpolation method to the minimum number of violated
constraints. Integer costs give an exact insertion identity, and a
nonnegative cubic remainder yields the local comparison. Finite binomial
replacement followed by Poisson averaging compares the endpoints without
differentiating an interpolation parameter. A three-point law of Boolean
warning rates and 128 moments reduce the quantitative claim to a finite
rational certificate. We formalize the complete argument in Lean 4 with
mathlib, including the analytic enclosures, the numerical inequality,
de-Poissonization, concentration, and transfer to distinct-variable clauses.
The formal theorem uses only the standard logical axioms of Lean; no
external numerical oracle is assumed.
\end{abstract}

\section{Introduction}
The random 3-SAT problem asks how the satisfiability of a random Boolean
formula depends on its number of clauses per variable. We give a complete
proof at the explicit density
\[
  \alpha=\frac{1067}{250}=4.268.
\]
The numerical value is accompanied by a finite rational certificate and a
Lean formalization of its implication for the random formula model.

For $n\geq3$, a legal clause is obtained by choosing an ordered triple of
distinct variables uniformly and independently choosing three fair
Boolean signs. A sign records the value falsifying the corresponding
literal. Clauses are sampled independently, so repetitions of an entire
clause are permitted. Ordering the variables changes no probability of
interest: every unordered signed clause has exactly six ordered
representations.

\begin{theorem}\label{thm:main}
Let $F_n$ consist of $\lfloor\alpha n\rfloor$ independent legal clauses
on $n$ Boolean variables, where $\alpha=1067/250$. Then
\[
  \lim_{n\to\infty}\Prob(F_n\text{ is satisfiable})=0.
\]
\end{theorem}

This statement neither assumes the existence of a limiting satisfiability
threshold nor identifies its exact value. It is a probabilistic statement
about random formulas, not an algorithm for deciding arbitrary instances.

\paragraph{Context and contribution.}
Franz and Leone~\cite{FL} introduced replica bounds for diluted systems
using a clause/site Poisson interpolation. Panchenko and
Talagrand~\cite{PT} developed a broader formulation of such bounds, and
Lelarge and Oulamara~\cite{LO} gave a combinatorial interpolation for
general degree distributions. The present proof follows this line of
comparison arguments. It uses integer minimum costs to obtain exact
insertion identities, and proves the required cubic sign for nonnegative
features directly. This qualification matters: the cited preprint of
Franz--Leone makes additional overlap assumptions in its odd-arity
discussion. We do not invoke that discussion as an unconditional 3-SAT
theorem.

Mertens, M\'ezard, and Zecchina~\cite{MMZ} studied threshold predictions
from the cavity equations. A distinct combinatorial upper-bound argument
of D\'iaz, Kirousis, Mitsche, and P\'erez-Gim\'enez~\cite{DKMP} proves
unsatisfiability at density $4.4898$. Our purpose here is to establish the
stated $4.268$ bound by a self-contained argument and an explicit,
kernel-checked calculation. We make no priority claim for the interpolation
method or a claim to have optimized the warning law.

\paragraph{Proof outline.}
We first sample the three variable positions with replacement. Write
$\en(F)$ for the minimum number of violated clauses. A comparison with
random unit constraints bounds the mean of $-y\en(F)$ in a Poisson clause
model. A rational certificate makes the bound negative for one fixed,
sufficiently large $y>0$. De-Poissonization then gives
$\E\en(F)\geq cn$ for some $c>0$. The bounded effect of one clause implies
$\Var(\en(F))\leq\alpha n$, so Chebyshev's inequality bounds the
satisfiability probability by $O(1/n)$. Conditioning every clause to have
distinct variables costs only a constant factor.

\section{Integer insertion and the comparison model}
\subsection{An exact insertion identity}
Let $\Omega$ be a nonempty finite set and let $H:\Omega\to\mathbb{N}$.
Write $\en(H)=\min_{\sigma\in\Omega}H(\sigma)$.

\begin{lemma}\label{lem:integer}
For every predicate $P$ on $\Omega$,
\begin{equation}\label{eq:integer}
 \en(H+\ind_P)=\en(H)+
 \ind\{P(\sigma)\text{ holds for every minimizer of }H\}.
\end{equation}
\end{lemma}
\begin{proof}
If some minimizer does not satisfy $P$, it still has cost $\en(H)$ after
insertion, and no cost decreases. Otherwise every old minimizer costs
$\en(H)+1$. Each other assignment already costs at least $\en(H)+1$,
because costs are integers. An old minimizer attains this new lower bound.
\end{proof}

For $y>0$ and $q=1-e^{-y}$, the exponential form of this identity is
\begin{equation}\label{eq:expinsert}
 e^{-y\en(H+\ind_P)}=e^{-y\en(H)}(1-qB_P),
\end{equation}
where $B_P$ is the indicator on the right-hand side of
\eqref{eq:integer}.

\subsection{Rows, fields, and their order of averaging}
Fix $n\geq1$. A seed $x=(i,J,R)$ consists of a uniform variable
$i\in\{1,\ldots,n\}$, a fair Boolean falsifying sign $J$, and an
independent rate $R$ with law
\begin{equation}\label{eq:law}
 \Prob(R=1/20)=1/4,\qquad
 \Prob(R=2/5)=1/2,\qquad
 \Prob(R=17/20)=1/4.
\end{equation}
A row consists of three independent seeds $x_1,x_2,x_3$.
All rows are independent. Given a row, sample independent Boolean
warnings $W_2,W_3$ with success probabilities $R_2,R_3$.
There are two kinds of constraint on $\sigma\in\{0,1\}^n$:
\begin{align*}
 C_x(\sigma)&=\ind\{\sigma_{i_1}=J_1,
                             \sigma_{i_2}=J_2,\sigma_{i_3}=J_3\},\\
 S_{x,W}(\sigma)&=\ind\{\sigma_{i_1}=J_1\}W_2W_3.
\end{align*}
The first is a clause with replacement in its variable positions; the
second is an active unit constraint when both warnings occur. Unused
coordinates let both constraints share the same row law.

For $K$ clauses and $L$ sites, hold all rows fixed and let $a$ denote the
array of site warnings. Let $H_a$ be the sum of the constraint costs and
$E_a=\en(H_a)$. Define
\begin{equation}\label{eq:functional}
 Z_y=\E_a e^{-yE_a},\qquad \Phi_y=\log Z_y,\qquad
 f(K,L)=\E_{\rm rows}\Phi_y.
\end{equation}
Thus the warning average is taken before the logarithm, and the row
average after it. Conditional on the rows, all warnings are independent.
All spaces at fixed $K,L$ are finite. Since $0\leq E_a\leq K+L$,
\begin{equation}\label{eq:growth}
 -y(K+L)\leq\Phi_y\leq0.
\end{equation}
Inserting one clause or site changes $\Phi_y$ by an amount in $[-y,0]$:
under each warning realization the minimum increases by either zero or
one, and averaging preserves the resulting exponential bounds.
The same increment bounds hold for $f$.

\section{The local cubic inequality}
At fixed existing rows, tilt the existing warning law by
\[
 \mu_y(a)=\frac{\Prob(a)e^{-yE_a}}{Z_y}.
\]
For a fresh seed $x=(i,J,R)$ define
\[
 b(x,a)=\ind\{\sigma_i=J\text{ for every minimizer of }H_a\},
 \qquad r(x,a)=R.
\]
Both take values in $[0,1]$. Let $\langle\cdot\rangle$ denote expectation
with respect to $\mu_y$.

A clause is violated by every minimizer exactly when each of its three
literals is falsified by every minimizer. This is the elementary
commutation of the two universal quantifiers over minimizers and literal
positions; it remains valid when variable positions repeat. Consequently,
\eqref{eq:expinsert} gives the clause increment
\begin{equation}\label{eq:clause}
 \Delta_C\Phi_y=\log\bigl(1-q\langle b(x_1,a)b(x_2,a)b(x_3,a)\rangle\bigr).
\end{equation}
For a site, first average its two fresh warnings, which are independent
of the existing warning array. Their product has mean $R_2R_3$, hence
\begin{equation}\label{eq:site}
 \Delta_S\Phi_y=\log\bigl(1-q\langle b(x_1,a)R_2R_3\rangle\bigr).
\end{equation}
Define the correction
\begin{equation}\label{eq:dy}
 d_y=\E\log(1-qR_1R_2R_3),
\end{equation}
where the rates are independent copies of~\eqref{eq:law}.

\begin{lemma}\label{lem:cubic}
Averaging the fresh row gives
\begin{equation}\label{eq:local}
 \E_x\Delta_C\Phi_y-3\E_x\Delta_S\Phi_y+2d_y\leq0.
\end{equation}
\end{lemma}
\begin{proof}
For $0\leq z\leq1$,
$\log(1-qz)=-\sum_{h\geq1}q^hz^h/h$. The absolute sum is bounded by
$\sum_{h\geq1}q^h/h=y$, uniformly in $z$, so it may be averaged termwise.
For a fixed $h$, sample independent replicas $a_1,\ldots,a_h$ from the
same law $\mu_y$, and set
\[
 B_h=\E_x\prod_{t=1}^h b(x,a_t),\qquad C_h=\E R^h.
\]
The seed $x$ is shared across the $h$ factors in $B_h$. Independence of
the three fresh seeds and finite Fubini give, respectively, the replica
averages of $B_h^3$, $B_hC_h^2$, and $C_h^3$ for the three moments in
\eqref{eq:local}. Their combination is nonnegative, because
\[
 B_h^3-3B_hC_h^2+2C_h^3=(B_h-C_h)^2(B_h+2C_h)\geq0.
\]
Multiplication by the nonpositive logarithmic coefficients proves the claim.
\end{proof}

Averaging the existing rows in~\eqref{eq:local} yields
\begin{equation}\label{eq:difference}
 f(K+1,L)-3f(K,L+1)+2f(K,L)+2d_y\leq0.
\end{equation}
The cubic sign here uses only nonnegative features; no overlap hypothesis
is needed. It is the cubic instance of the polynomial remainder in
Lelarge--Oulamara~\cite[condition (4), equation (30)]{LO}.

\section{Finite replacement and Poisson endpoints}
The next comparison is elementary and independent of the SAT model.

\begin{lemma}\label{lem:comparison}
Suppose $f:\mathbb{N}^2\to\mathbb{R}$ has uniformly bounded increments
in each coordinate and satisfies~\eqref{eq:difference} with a constant
$d$ in place of $d_y$. For every $a\geq0$,
\begin{equation}\label{eq:comparison}
 \E_{K\sim\Pois(a)}f(K,0)
 \leq\E_{L\sim\Pois(3a)}f(0,L)-2ad.
\end{equation}
\end{lemma}
\begin{proof}
Define
\[
 T_N(K,L)=\sum_{j=0}^N\binom Nj2^{N-j}f(K+j,L).
\]
Pascal's identity gives
$T_{N+1}(K,L)=T_N(K+1,L)+2T_N(K,L)$. Induction implies
\begin{equation}\label{eq:binomial}
 T_N(K,L)\leq3^N\bigl(f(K,L+N)-2Nd/3\bigr).
\end{equation}
Indeed, the case $N=0$ is equality. Applying the induction hypothesis to
both summands of $T_{N+1}$ gives an upper bound
\[
 3^N\bigl(f(K+1,L+N)+2f(K,L+N)-2Nd\bigr).
\]
The assumed local inequality at $(K,L+N)$ bounds this by
$3^{N+1}(f(K,L+N+1)-2(N+1)d/3)$.

At $K=L=0$, divide~\eqref{eq:binomial} by $3^N$ and average
$N\sim\Pois(3a)$. Conditional on $N$, the left side averages
$f(J,0)$ with $J\sim\Bin(N,1/3)$. For $N=j+\ell$, its weight is
\[
 e^{-3a}\frac{(3a)^N}{N!}\binom Nj(1/3)^j(2/3)^\ell
 =e^{-3a}\frac{a^j}{j!}\frac{(2a)^\ell}{\ell!}.
\]
Summing $\ell$ gives the $\Pois(a)$ weight of $j$, whereas
$\E N=3a$. Bounded increments give at most linear growth, justifying
all rearrangements by absolute convergence. The case $a=0$ is immediate.
\end{proof}

Let $F^{\rm rep}_K$ have $K$ clauses whose variable positions are sampled
with replacement. At the clause-only endpoint the inner warning average
disappears, so $f(K,0)=-y\E\en(F^{\rm rep}_K)$. With $a=\alpha n$,
Lemma~\ref{lem:comparison} becomes
\begin{equation}\label{eq:poissonenergy}
 -\frac yn\E\en(F^{\rm rep}_{\Pois(\alpha n)})
 \leq\frac1n\E_{L\sim\Pois(3\alpha n)}f(0,L)-2\alpha d_y.
\end{equation}
Here a Poisson subscript means that the count is sampled independently
and then the indicated number of independent clauses is sampled.

\section{The site endpoint and its moments}
\subsection{Conflicting warnings}
Hold the site rows fixed. For each variable $v$ and sign $s\in\{0,1\}$,
let $N_{v,s}$ count sites rooted at $(v,s)$ whose two warnings are active.
Every assignment violates at least one site at a variable where both
$N_{v,0}$ and $N_{v,1}$ are positive. Therefore the minimum site energy
is at least the number of such variables.

Put
\[
 A_v=\prod_{\text{sites rooted at }(v,0)}(1-R_2R_3),\quad
 B_v=\prod_{\text{sites rooted at }(v,1)}(1-R_2R_3),\quad
 G_v=A_v+B_v-A_vB_v.
\]
An empty product is one. The two sign groups use disjoint warning
coordinates, so $G_v$ is the conditional probability that at least one
of their active-warning counts is zero. Different variables also use
disjoint coordinates. Writing $\delta=e^{-y}$, the energy lower bound and
this conditional independence imply
\begin{equation}\label{eq:vertex}
 \Phi_y\leq\sum_{v=1}^n\log\bigl(G_v+(1-G_v)\delta\bigr).
\end{equation}
This is an inequality: counting conflicting variables need not give the
exact minimum energy.

\begin{lemma}\label{lem:truncate}
For $x\in[0,1]$, $\delta\in(0,1]$, and an integer $T\geq1$,
\begin{equation}\label{eq:truncate}
 \log(1-(1-\delta)x)\leq-\sum_{k=1}^T\frac{x^k}{k}+T\delta.
\end{equation}
\end{lemma}
\begin{proof}
Truncating the logarithmic series gives an upper bound
$-\sum_{k=1}^T(1-\delta)^kx^k/k$. Bernoulli's inequality gives
$(1-\delta)^k\geq1-k\delta$, and $x^k\leq1$ bounds the total correction
by $T\delta$. This also covers $x=1$ and $\delta=1$.
\end{proof}

\subsection{Mixed Poisson moments}
Let $L\sim\Pois(3\alpha n)$ and fix a variable $v$. Define
\begin{equation}\label{eq:moments}
 \nu_j=\E(1-RR')^j,\qquad
 M_j=\exp\bigl((3\alpha/2)(\nu_j-1)\bigr),\qquad
 \mu_k=\sum_{j=0}^k(-1)^j\binom kj M_j.
\end{equation}
All rate copies in the same expectation are independent. For independent
rows $X_i$ and a function $0\leq h\leq1$, conditioning on $L$ gives
\[
 \E\prod_{i=1}^L h(X_i)
 =\exp\bigl(3\alpha n(\E h(X)-1)\bigr).
\]
A row is rooted at either given sign of $v$ with probability $1/(2n)$,
and cannot belong to both sign groups. Its contribution to $A_v^jB_v^\ell$
therefore has mean
$1+(\nu_j-1)/(2n)+(\nu_\ell-1)/(2n)$. Hence
\begin{equation}\label{eq:mixed}
 \E[A_v^jB_v^\ell]
 =\exp\bigl((3\alpha/2)(\nu_j+\nu_\ell-2)\bigr)=M_jM_\ell.
\end{equation}
Finite binomial expansion now yields
\begin{equation}\label{eq:musquare}
 \E\bigl[((1-A_v)(1-B_v))^k\bigr]=\mu_k^2.
\end{equation}
In particular $\mu_k=\E(1-A_v)^k\geq0$.

Apply Lemma~\ref{lem:truncate} with $x=(1-A_v)(1-B_v)=1-G_v$,
then average~\eqref{eq:vertex} and use~\eqref{eq:musquare}. We obtain
\begin{equation}\label{eq:sitebound}
 \frac1n\E_{L\sim\Pois(3\alpha n)} f(0,L)
 \leq-\sum_{k=1}^T\frac{\mu_k^2}{k}+Te^{-y}.
\end{equation}
Only a finite polynomial has been integrated. In particular, no
integrability claim for $\log G_v$ or $1/G_v$ is needed.

Since $q\leq1$ and $R_1R_2R_3<1$, equation~\eqref{eq:dy} implies
$d_y\geq\E\log(1-R_1R_2R_3)$. Define
\begin{equation}\label{eq:psi}
 \Psi_T=-\sum_{k=1}^T\frac{\mu_k^2}{k}
             -2\alpha\E\log(1-R_1R_2R_3).
\end{equation}
Combining~\eqref{eq:poissonenergy} and~\eqref{eq:sitebound} proves
\begin{proposition}\label{prop:energy}
For every $n\geq1$, $y>0$, and $T\geq1$,
\begin{equation}\label{eq:energybound}
 -\frac yn\E\en(F^{\rm rep}_{\Pois(\alpha n)})\leq\Psi_T+Te^{-y}.
\end{equation}
\end{proposition}

\section{The rational certificate}\label{sec:certificate}
For the law~\eqref{eq:law} and $T=128$, the numerical statement is
\begin{equation}\label{eq:negative}
 \Psi_{128}\leq-\frac1{100000}=-\varepsilon.
\end{equation}
We specify a rational calculation proving this inequality. The displayed
decimal approximations below are explanatory only.

Represent the rate law by four equally likely labels
\[
 (r_0,r_1,r_2,r_3)=(1/20,17/20,2/5,2/5).
\]
Then $\nu_j=\frac1{16}\sum_{i,\ell=0}^3(1-r_ir_\ell)^j$ is rational.
Put $x_j=(3\alpha/2)(\nu_j-1)$ for $0\leq j\leq128$.

\subsection{Exponential enclosures}
For an integer scale $s>0$, define
$D_s(z)=\lfloor sz\rfloor/s$ and $U_s(z)=-D_s(-z)$. Define intervals
recursively by
\begin{align*}
 I_0(x)&=[\max(0,1+x),\,1/(1-x)],\\
 I_{t+1}(x)&=[D_s(l^2),U_s(u^2)],
       \qquad [l,u]=I_t(x/2).
\end{align*}
For $x<2^t$, the interval $I_t(x)$ encloses $e^x$ and has nonnegative
endpoints. At $t=0$, use $1+x\leq e^x$ and
$1-x\leq e^{-x}$; the latter can be inverted since $x<1$.
Induction uses $e^x=(e^{x/2})^2$ and outward rounding.

Compute $I_{160}(x_j)$ on the grid $s=2^{320}$, then round its endpoints
outward on the grid $2^{160}$ to obtain $[l_j,u_j]$. The domain condition
holds because $x_j\leq0$. For $c_{kj}=(-1)^j\binom kj$, put
\[
 L_k=\sum_{j=0}^k c_{kj}
 \begin{cases}l_j,&c_{kj}\geq0,\\u_j,&c_{kj}<0,\end{cases}
 \qquad V=-\sum_{k=1}^{128}\frac{\max(0,L_k)^2}{k}.
\]
Signed interval arithmetic gives $L_k\leq\mu_k$. Since $\mu_k\geq0$,
$\max(0,L_k)^2\leq\mu_k^2$, so $V$ bounds the vertex term from above.
The clamp before squaring prevents a negative interval endpoint from
giving an invalid lower bound on a square.

\subsection{Logarithmic enclosures and exact arithmetic}
For $0\leq z<1$, let
\[
 \lambda(z)=-\sum_{h=1}^{32}\frac{z^h}{h}-\frac{z^{33}}{1-z}.
\]
Bounding the remaining logarithmic-series tail by a geometric series
proves $\lambda(z)\leq\log(1-z)$. Thus
\[
 W=-\frac{2\alpha}{64}\sum_{i,j,\ell=0}^3
                   \lambda(r_ir_jr_\ell)
\]
is an upper bound on the edge term. The resulting entirely
rational inequality is
\begin{equation}\label{eq:rational}
 V+W\leq-1/100000.
\end{equation}
Lean computes and proves~\eqref{eq:rational} by \code{decide +kernel},
and the analytic enclosure lemmas then imply~\eqref{eq:negative}.
The exponential intervals are computed within Lean from their recurrence;
no table supplied by an external program is a premise of the theorem.

An independent Python implementation uses rational Taylor enclosures
for the exponential and a different logarithm enclosure. Its upper
bounds, rounded for display, are
\begin{center}
\begin{tabular}{lr}
\toprule
Contribution & Independent upper bound \\
\midrule
Vertex & $-0.7611927541550103$\\
Edge & $+0.7611815516097719$\\
Sum & $-0.000011202545238429759$\\
\bottomrule
\end{tabular}
\end{center}
These numbers cross-check the calculation; they are not inputs to Lean
and are not claimed to be exact values of the real expression $\Psi_{128}$.

\section{Fixed clause count, concentration, and legal clauses}
Choose one finite $y>0$ such that $128e^{-y}\leq\varepsilon/4$, and keep
it fixed as $n$ grows. Let $M=\lfloor\alpha n\rfloor$ and
$N\sim\Pois(\alpha n)$. Coupling formulas by prefixes of the same
independent clause sequence and using the one-clause energy bound gives
\begin{align}
 \bigl|\E\en(F^{\rm rep}_N)-\E\en(F^{\rm rep}_M)\bigr|
 &\leq\E|N-M|\nonumber\\
 &\leq\sqrt{\E(N-M)^2}\leq\sqrt{\alpha n+1}.\label{eq:depois}
\end{align}
Here $\E(N-M)^2=\alpha n+(\alpha n-M)^2$. Equations
\eqref{eq:energybound}, \eqref{eq:negative}, and~\eqref{eq:depois} yield
\[
 -\frac yn\E\en(F^{\rm rep}_M)
 \leq-\varepsilon+128e^{-y}+\frac yn\sqrt{\alpha n+1}.
\]
Eventually the last term is at most $\varepsilon/4$, so, for
$c=\varepsilon/(2y)>0$,
\begin{equation}\label{eq:linear}
 \E\en(F^{\rm rep}_M)\geq cn
\end{equation}
for all sufficiently large $n$. This fixes $y$ before taking $n\to\infty$;
there is no exchange of two limits.

Replacing one clause changes every assignment's violation count by at
most one, hence changes its minimum by at most one. To bound the variance,
reveal the $M$ clauses one at a time and use the Doob martingale of
$\en(F^{\rm rep}_M)$. Conditional on the preceding clauses, the
conditional expectations corresponding to different values of the next
clause differ by at most one: couple all unrevealed clauses identically.
Each martingale increment therefore has absolute value at most one.
Its conditional mean is zero, so orthogonality gives
\begin{equation}\label{eq:variance}
 \Var(\en(F^{\rm rep}_M))\leq M\leq\alpha n.
\end{equation}
The formal proof uses the equivalent finite-product variance decomposition.
Chebyshev's inequality and~\eqref{eq:linear} imply eventually
\begin{equation}\label{eq:chebyshev}
 \Prob(\en(F^{\rm rep}_M)=0)\leq\frac{\alpha}{c^2n}.
\end{equation}

Let $A$ be the event that each of the $M$ clauses has distinct variable
positions. Conditional on $A$, clauses are independent uniform legal
clauses and their signs remain independent and fair. For $n\geq6$,
\begin{equation}\label{eq:conditioning}
 \Prob(A)^{-1}=
 \left(\frac{n^2}{(n-1)(n-2)}\right)^M
 \leq e^{6M/n}\leq e^{6\alpha}.
\end{equation}
Indeed, for $r=n^2/((n-1)(n-2))$, use
$\log r\leq r-1\leq6/n$; the latter inequality reduces to
$3n^2-16n+12\geq0$, valid for $n\geq6$.
Since $\Prob(B\mid A)\leq\Prob(B)/\Prob(A)$ for any event $B$,
\eqref{eq:chebyshev} and~\eqref{eq:conditioning} give
\[
 \Prob(F_n\text{ is satisfiable})
 \leq e^{6\alpha}\frac{\alpha}{c^2n}\longrightarrow0.
\]
This proves Theorem~\ref{thm:main}. The conditioning constant is deliberately
coarse; the argument supplies an asymptotic conclusion, not a useful
small-instance estimate.

\section{Relation to the partition-moment formulation}
For the same finite Hamiltonian $H_a$ and field law on both sides, put
$Z_\beta(a)=\sum_\sigma e^{-\beta H_a(\sigma)}$. For $\beta\geq0$,
\[
 e^{-\beta E_a}\leq Z_\beta(a)\leq2^n e^{-\beta E_a}.
\]
Raise to a power $m\geq0$, average over the fields, and take logarithms:
\begin{equation}\label{eq:bridge}
 \Phi_{m\beta}\leq\log\E_a[Z_\beta(a)^m]
                  \leq\Phi_{m\beta}+mn\log2.
\end{equation}
Thus at fixed $m\beta=y$ the normalized difference is at most $m\log2$.
This finite-system sandwich, also formalized in Lean, relates our
functional to a nested partition moment of the form in
Franz--Leone~\cite[equation (45)]{FL}.

The comparison in Section~4 retains the clause/site balance of their
Poisson interpolation~\cite[equations (51)--(54)]{FL}, but proves it
using finite replacement. Equation~\eqref{eq:bridge} uses an identical
fixed field law on its two sides; it does not identify that law with the
reweighted measure in their equation (46). Nor do we claim the exact
integer-cost implementation appears in that paper. The direct insertion
identity is sufficient here, at the cost of restricting the argument to
integer unit costs rather than a general finite-temperature model.

\section{Formalization and reproducibility}\label{sec:lean}
The formalization uses Lean 4.24.0 and mathlib~\cite{mathlib}, pinned to
the revision
\begin{center}
\code{f897ebcf72cd16f89ab4577d0c826cd14afaafc7}.
\end{center}
The public
model expresses probability as the ratio of two finite cardinalities.
Clauses are triples with an injective variable map and three Boolean
falsifying signs; a formula is an ordered list of independent clauses.
The exact expanded endpoint is
\begin{verbatim}
Filter.Tendsto
  (fun n : Nat => SatUpper.satisfiabilityProbability n (1067*n/250))
  Filter.atTop (nhds 0)
\end{verbatim}
Natural-number division gives exactly $\lfloor1067n/250\rfloor$.
For $n<3$, the legal clause set is empty and the definition uses Lean's
division-by-zero convention; these finitely many indices do not affect
the limit. The final theorem has no undischarged interpolation, numerical,
or concentration hypotheses.

\begin{center}
\small
\begin{tabular}{>{\raggedright\arraybackslash}p{.42\textwidth}>{\raggedright\arraybackslash}p{.50\textwidth}}
\toprule
Mathematical step & Modules below \code{lean/SatUpper/}\\
\midrule
Model and final theorem & \code{Model}, \code{Threshold}\\
Integer insertion & \code{Ground/Minimum}, \code{Ground/Finite}\\
Insertion and cubic sign & \code{Ground/Insertion}, \code{LogInsertion}, \code{ReplicaCubic}\\
Finite replacement and Poisson average & \code{BinomialPoisson}, \code{PoissonExpectation}, \code{Ground/Comparison}\\
Site bound and mixed moments & \code{Ground/Site}, \code{Ground/SiteBound}, \code{PoissonPair}\\
Rational certificate & \code{GridExponential}, \code{NumericalMoments}, \code{NumericalLogs}, \code{NumericalEnclosure}\\
Fixed count and asymptotics & \code{ReplacementPoissonMean}, \code{ModelTransfer}, \code{Obligations}\\
Variance and conditioning & \code{FiniteVariance}, \code{VarianceControl}, \code{LegalConditioning}\\
Partition-moment sandwich & \code{Ground/Minimum}, \code{Ground/Finite}\\
\bottomrule
\end{tabular}
\end{center}

The audit enumerates the project theorems, asks Lean for their transitive
axiom dependencies, and separately checks the expanded endpoint above.
It allows only \code{propext}, \code{Classical.choice}, and
\code{Quot.sound}. It rejects proof placeholders and additional
evaluation axioms. This is a dependency audit using Lean's environment,
not a separate proof checker: mathematical checking is performed by the
Lean kernel. The definitions of the model and the correspondence to the
paper still require human inspection.

With the pinned toolchain installed, the proof and audit are reproduced
from the repository root by
\begin{verbatim}
sh lean/run.sh exe cache get
python3 lean/verify.py
python3 lean/test_verify.py
\end{verbatim}
The independent certificate uses only the Python standard library:
\begin{verbatim}
python3 src/richer_certificate.py --alpha 4.268 --cutoff 128 \
  --surveys .05,.85,.5 .4,.4,.5 --output certificate.json
\end{verbatim}
The repository additionally includes numerical cross-checks, source
provenance, and instructions for rebuilding the paper and release bundle.
\paragraph{Availability.}
The Lean sources, independent certificate code, and reproduction instructions
are available at
\url{https://github.com/yamelton/sat-upper-4268}.
The versioned software archive for release \texttt{v1.0.0} is preserved at
\url{https://doi.org/10.5281/zenodo.23120101}.
The software is licensed under Apache-2.0, with
third-party notices preserved. The manuscript is excluded from the software
license; copyright in the manuscript is retained by Fedor Vorobyev.


\paragraph{Acknowledgments and tools.}
The development relies on the Lean and mathlib communities and credits
the interpolation methods cited above. AI assistance was used in preparing
the manuscript and publication materials. The author is responsible for
the mathematical claims and the final text.

\bibliographystyle{plain}
\bibliography{references}
\end{document}
